\subsection{More about F-term uplifting --- arXiv:0707.2671}

\textbf{Authors:} Hiroyuki Abe, Tetsutaro Higaki, Tatsuo Kobayashi

\paragraph{主な主張}
uplifting場とK\"ahlerモジュラスが超ポテンシャルで直接混合する模型でもF-term upliftingが実現し、重いモジュラスとmirage mediation型gaugino質量が両立することを示した\cite{Abe:2007yb}。

\paragraph{新規性}
非摂動的な場の混合によってモジュラス質量とgravitino質量の階層が増大しても、モジュラスF項の抑制は同じ比率では強まらないことを明らかにした。

\paragraph{位置付け}
弦インスタントン由来の結合を取り入れ、SUSY破れセクターとモジュラスを分離する近似を拡張した研究。

\paragraph{コメント（読書メモ）}
檜垣さんに挙げていただいた論文。$W\sim e^{-bT}X$型の結合を持つF-term upliftingとして参照。
